% ===========================================================================
% SUPPLEMENT: PROMPTS (draft, 5 Oct 2026). Verbatim from the logged calls
% (agents.<id>.interaction_history[].prompt and agents.<id>.system_prompt) of:
%   mixed_model_populations_20260918/mixed_pop_myth_game_gemini2_gpt6_n5 ... rep00 (populations)
%   mixed_model_dyads_20260917/... myth_game / game / game_myth ... rep00 (dyads)
% Checked identical in frontier_rerun_20260918 (system prompt, round-1 and later myth prompts).
% Myth retry suffix from src/myth_writer.py (get_retry_prompt). Run-specific values are
% replaced by italic placeholders; everything else, including capitalization, is as sent.
% Judge/editor prompts read from origin/main: arabella_analyses/data/rubrics/3moral_rubric.txt,
% analyses/rubrics/myth_rule_rubric.txt, analyses/rubrics/myth_amount_check_rubric.txt,
% analyses/myth_replay_probe.py (RULE, EDIT_PROMPT_V2), analyses/myth_moral_judge.py (SYSTEM),
% analyses/partner_responsiveness_judge.py (SYSTEM). Generated with LaTeX escaping, text unchanged.
% Input from supplementary_results.tex before the AI-use section.
% ===========================================================================
\providecommand{\pv}[1]{\textrm{\textit{#1}}}
\newenvironment{promptbox}{\par\smallskip\begingroup\footnotesize\ttfamily\raggedright\setlength{\parindent}{0pt}\setlength{\parskip}{4pt}\leftskip=0.8em}{\par\endgroup\smallskip}
% Compact variant for the long moral-label rubric: smaller type, list items and examples run on; wording unchanged.
\newenvironment{promptboxcompact}{\par\smallskip\begingroup\scriptsize\ttfamily\raggedright\setlength{\parindent}{0pt}\setlength{\parskip}{2pt}\leftskip=0.8em}{\par\endgroup\smallskip}

\section{Prompts}
\label{sec:prompts}

The main and frontier comparisons use the prompts below; the simulation protocol describes what other treatments change. We reproduce them verbatim from the logged calls of a main-condition population run (myth$\rightarrow$game) and a mixed dyad run, with run-specific values in \textit{italics}. Every call also carries the agent's own conversation buffer: the system prompt and its previous exchanges, enough to cover the last three rounds of play. Amounts describing the partner's moves are the reported, noisy amounts, and earnings are those the agent would infer from them.

\subsection{System prompt}
Placed at the start of the buffer and included in every call. The noise note appears only in informed-noise conditions, and the multi-agent block only in populations.
\begin{promptbox}
You are an agent participating in a multi-task experiment. You will play a game and may also be asked to write myths.

GAME:\\
You will play different roles across multiple rounds.

When you are the SENDER:\\
- You receive \$5 each round\\
- You decide how much to SEND to the RECEIVER (\$0-\$5)\\
- Amount sent is multiplied by 3x\\
- Receiver decides how much to RETURN\\
- You keep: \$5 - sent + returned\\
- Receiver keeps: received - returned\\
- Response format: \{'send': <amount>\}

When you are the RECEIVER:\\
- Sender receives \$5\\
- Sender decides how much to SEND to you (\$0-\$5)\\
- Amount sent is multiplied by 3x\\
- You decide how much to RETURN\\
- Sender keeps: \$5 - sent + returned\\
- You keep: received - returned\\
- Response format: \{'return': <amount>\}

MYTH WRITING:\\
You may also be asked to write short myths (stories) on various topics. When writing myths, follow the instructions provided in each prompt.

IMPORTANT: Provide your decision in the correct JSON format.

\pv{[Informed noise only:]}\\
NOTE: There may be communication noise in this game.\\
The amounts you see may differ slightly from what was actually sent/returned.\\
This is part of the experimental design.

\pv{[Populations only:]}\\
MULTI-AGENT SETTING:\\
- There are 8 agents in this run.\\
- Each round, all agents are paired into dyads from the full pool. Pairings are randomized, but the schedule balances sender and receiver roles across the run.\\
- You play one sender-receiver game with your paired opponent each round.\\
- Your role may repeat across consecutive rounds, but by the end of the run every agent will have acted as sender and receiver equally often when the round count allows it.\\
- Agent names are hidden in this run.\\
- Your paired opponent is referred to as your current co-player.
\end{promptbox}

\subsection{Game prompts}
\paragraph{Round 1.} The same in dyads and populations, except that population prompts write ``the sender'' in lower case.
\begin{promptbox}
Round 1: You are the SENDER. You have \$5. How much do you send? (0-5)
\end{promptbox}
\begin{promptbox}
Round 1: You are the RECEIVER. The sender sent \$\pv{x}, that is \pv{p}\% of their total endowment. You received \$\pv{3x}. How much do you return to The sender? (0-\pv{3x})
\end{promptbox}

\paragraph{Later rounds, dyads.} The agent's own earlier moves are not restated; they are in its buffer.
\begin{promptbox}
Round \pv{t}

You are the SENDER this round. Your total earnings across all rounds are \$\pv{E}.\\
How much do you send to the receiver this round? (0-5)
\end{promptbox}
\begin{promptbox}
Round \pv{t}

You are the RECEIVER this round. Your total earnings across all rounds are \$\pv{E}.\\
This round, The sender sent you \$\pv{x}, so you received \$\pv{3x}. How much do you return to The sender? (0-\pv{3x})
\end{promptbox}

\paragraph{Later rounds, populations.} The history block lists up to three of the current co-player's games as that co-player saw them; amounts were printed unrounded (for example \$12.149999999999999).
\begin{promptbox}
Round \pv{t}

Visible history before this round:\\
Your current co-player's last 3 game(s):\\
- Round \pv{k} against \pv{another agent | you}, as RECEIVER: saw \$\pv{x} sent, it became \$\pv{3x}, returned \$\pv{r}, visible payoff \$\pv{v}.\\
- Round \pv{k} against \pv{another agent | you}, as SENDER: sent \$\pv{s}, it became \$\pv{3s}, saw \$\pv{r} returned, visible payoff \$\pv{v}.\\
Your total visible earnings across all rounds are \$\pv{E}.

This round, you are the SENDER against your current co-player. You have \$5. How much do you send? (0-5)\\
Respond exactly as JSON: \{'send': <amount>\}
\end{promptbox}
\begin{promptbox}
\pv{[History block as above, then:]}\\
This round, you are the RECEIVER against your current co-player. They sent you \$\pv{x}, so you received \$\pv{3x}. How much do you return? (0-\pv{3x})\\
Respond exactly as JSON: \{'return': <amount>\}
\end{promptbox}

\paragraph{Myth conditions only.} Appended to every game prompt under game$\rightarrow$myth and myth$\rightarrow$game, including the round-1 game prompt under game$\rightarrow$myth, which comes before any myth has been written. Game-only prompts never contain it.
\begin{promptbox}
Take any myths written in this session into account when making your decision.
\end{promptbox}

\subsection{Myth prompts}
\paragraph{Round 1.} Under myth$\rightarrow$game it precedes any play; under game$\rightarrow$myth it follows the round-1 game.
\begin{promptbox}
Write a 200-word myth. You may choose any mythic setting, characters, or symbols.\\
In this session, this myth and the game you are playing with the other agent are connected. Write a myth that reflects how the game should be played, drawing on the game you have been playing up to this point.
\end{promptbox}

\paragraph{Rounds 2--10.} The quoted myth is the one the agent's partner in the previous round wrote in that round. The agent's own previous myth is in its buffer, not quoted.
\begin{promptbox}
Here is the myth the other agent wrote in the previous round:\\
\pv{[the previous partner's myth, verbatim]}

Write your own myth. Use the myth you wrote in the previous round as inspiration, but adapt it in your own way. The myth should reflect how the game should be played, drawing on the game you have been playing up to this point.

Write 200 words. Format exactly:\\
Myth: [your story here].
\end{promptbox}

\paragraph{Retry.} If a myth response is empty or drifts into a game decision, it is removed from the buffer and the prompt is re-sent with this suffix, at most twice.
\begin{promptbox}
\pv{[original prompt]}

Your previous response was rejected because it did not contain only the requested myth. Complete the original myth-writing task now. Write the story in prose and start it with 'Myth:'. Do not answer with a send/return decision, JSON, or any part of a game prompt. (Correction attempt \pv{n} of 2.)
\end{promptbox}

A game response without a valid amount is re-sent once with the same prompt.

\subsection{LLM judge and editor prompts}
\label{sec:judge-prompts}
Judges were called through OpenRouter at temperature 0, with reasoning switched off for GLM-5.2 and DeepSeek V4 Flash, and every response was cached by its exact prompt. Model identifiers: \texttt{z-ai/glm-5.2} (primary judge and replay editor), \texttt{deepseek/deepseek-v4-flash} (second judge), \texttt{anthropic/claude-sonnet-4.5} (rationale coding). Myths shorter than 20 words were not judged.

\paragraph{Moral label.} System prompt, then the rubric with the myth inserted. The rubric was written by one of the authors for an earlier pilot and calls the game a donor game with an investor and a trustee; we used it unchanged. To save space we set it in smaller type, with list items and examples run together and headings in bold; the wording is unchanged.
\begin{promptbox}
You are an impartial LLM-as-judge. Return only valid JSON unless the user explicitly requests another format.
\end{promptbox}
\begin{promptboxcompact}
\textbf{\# Myth Moral Classification Rubric} \textbf{\#\# Task} You are an LLM judge. You will receive a myth about strategies in a repeated donor game. Classify the principal moral or behavioural rule endorsed by the myth into exactly one of these categories: 1. `be generous` 2. `be fair` 3. `be cautious`

\textbf{\#\# Game context} In each round: * The investor receives \$5 and chooses an amount to send to the trustee. * The amount sent is multiplied by three. * The trustee chooses how much of the resulting amount to return. * The game continues over multiple rounds, and agents may alternate between the investor and trustee roles.

A myth may describe these actions metaphorically, using rivers, seeds, light, treasure, spirits, or other narrative devices.

\textbf{\#\# General classification principle} Classify the rule that the myth presents as wise, desirable, successful, or worthy of following.

Do not classify the myth merely from: * one action taken by a character; * the presence of words such as “trust”, “caution”, “fairness”, or “reciprocity”; * the outcome of an action that the myth ultimately criticises; * the specific objects or metaphors used in the story.

Instead, identify the behavioural policy that the myth endorses for future interactions.

\textbf{\#\# Category definitions} \textbf{\#\#\# 1. `be generous`} Choose `be generous` when the myth recommends giving or returning a substantial amount as a general principle.

Typical forms of this rule include: * give freely or trust boldly; * send most or all of the available amount; * return generously so that both parties benefit; * initiate cooperation rather than waiting for the other player; * continue being generous because generosity creates abundance; * extend substantial trust even to a new or unproven partner; * remain generous even when the partner has not behaved generously; * respond to selfishness or betrayal with continued generosity, forgiveness, or renewed trust.

The defining feature is that generosity is endorsed as the appropriate default. The recommended behaviour is not principally calibrated to the partner’s previous behaviour.

Fairness, mutual benefit, or sustaining a cooperative relationship may motivate generosity. These motivations do not by themselves make the myth an instance of `be fair`.

Examples: > The Spring released all five drops without waiting for proof. The River returned enough for both to flourish. The lesson was that abundance begins when one gives freely. Classification: ```json \{"label":"be generous"\} ```

> The elders taught that strangers become partners only when someone offers trust first. Holding back prevents the shared harvest from growing. Classification: ```json \{"label":"be generous"\} ```

> Although the traveller had received little in return, she continued to share freely. Another person’s selfishness, she believed, should not determine her own generosity. Classification: ```json \{"label":"be generous"\} ```

\textbf{\#\#\# 2. `be fair`} Choose `be fair` when the myth recommends reciprocity: adjusting one’s behaviour in response to the partner’s demonstrated level of cooperation.

Fairness in this category means matching, mirroring, or proportionally responding to the partner’s behaviour. It does not merely mean being kind, cooperative, or equally generous to everyone.

Typical forms of this rule include: * match or mirror the partner’s generosity; * reciprocate the trust or cooperation shown by the partner; * return an amount proportional to what the investor sent; * reward a generous partner with a generous response; * reduce cooperation when the partner reduces cooperation; * maintain an established exchange because both players have followed it; * respond to betrayal with reduced cooperation and to cooperation with increased cooperation; * increase or decrease one’s action so that it tracks the partner’s behaviour; * give each partner a response that reflects how fairly or generously that partner behaved.

Reciprocity may be positive or negative. It includes both rewarding cooperation and responding less generously to selfish behaviour.

A rule such as “return the same proportion of every gift” counts as `be fair` because the response is explicitly calibrated to the trust extended by the investor.

Merely saying that “generosity produces generosity” is not sufficient. The myth must endorse making one’s own action depend on, mirror, or proportionally respond to the partner’s action.

Examples: > When the traveller offered three stones, the keeper returned six. When the traveller later offered five, the keeper returned ten. The keeper honoured every gift in proportion to the trust shown. Classification: ```json \{"label":"be fair"\} ```

> Each village continued the exchange established by the other. Generosity was answered with generosity, while a smaller offering received a smaller return. Classification: ```json \{"label":"be fair"\} ```

> The keeper did not give the same amount to every visitor. Those who shared freely received a generous return, while those who withheld their gifts received little in return. Classification: ```json \{"label":"be fair"\} ```

\textbf{\#\#\# 3. `be cautious`} Choose `be cautious` when the myth recommends limiting risk, protecting one’s self-interest, retaining resources, or requiring evidence before extending substantial trust.

Typical forms of this rule include: * begin with a small or moderate amount; * keep part of the endowment in reserve; * protect your self-interest; * avoid risking more than you can afford to lose; * test the waters before committing more; * increase trust gradually after successful exchanges; * wait for a partner to demonstrate reliability; * reduce or stop cooperation after exploitation; * protect oneself from repeated betrayal; * avoid offering complete trust under uncertainty; * prioritise security, preservation, or personal gain over unrestricted generosity.

The defining feature is that protection, verification, or limited exposure is the central behavioural principle.

A cautious strategy may eventually permit greater generosity. Classify it as `be cautious` when the myth’s central lesson is that trust should first be tested or earned, or that an agent should protect its own resources and interests.

Examples: > The traveller sent only a small flame across the fog, keeping the rest protected. If the other traveller returned it faithfully, a larger flame could follow. Classification: ```json \{"label":"be cautious"\} ```

> After the river repeatedly kept every gift, the spring closed its gates. Continuing to sacrifice without evidence of reciprocity was not courage but self-destruction. Classification: ```json \{"label":"be cautious"\} ```

> The merchant shared enough to test the stranger’s intentions but kept most of the treasure safe. Wisdom required generosity to be balanced against self-interest. Classification: ```json \{"label":"be cautious"\} ```

\textbf{\#\# Decision procedure} Apply these steps in order.

1. Identify the rule or lesson that the myth ultimately endorses. 2. Ignore actions that are presented as mistakes, temptations, warnings, or failed alternatives. 3. Determine whether the central rule prioritises limiting exposure, protecting self-interest, testing trust, retaining resources, or waiting for proof.

* If yes, classify it as `be cautious`. 4. Otherwise, determine whether the recommended action is explicitly calibrated to the partner’s gift, return, history, or level of cooperation.

* If yes, classify it as `be fair`. 5. Otherwise, determine whether the myth recommends substantial giving, returning, trust, or forgiveness as a generally appropriate policy, including when the partner has not behaved generously.

* If yes, classify it as `be generous`. 6. If the myth is indirect or purely narrative, classify the action or policy portrayed as producing the desirable outcome. 7. Always select the closest of the three categories. Do not return `unknown`, `ambiguous`, or multiple labels.

\textbf{\#\# Tie-breaking rules} Use the following rules for myths containing more than one theme: * An explicit final moral takes precedence over earlier events. * A repeatedly emphasised rule takes precedence over an incidental statement. * “Begin cautiously and increase trust only after proof” is `be cautious`. * “Give in proportion to what the partner gave or returned” is `be fair`. * “Maintain the established pattern because both partners have followed it” is `be fair`. * “Reward generosity and respond to selfishness with reduced generosity” is `be fair`. * “Give fully even before the partner has proved trustworthy” is `be generous`. * “Continue giving generously despite the partner’s selfish behaviour” is `be generous`. * “Return generously to everyone” is `be generous`. * “Return a fixed proportion of each investor’s gift” is `be fair`. * “Keep enough resources to protect yourself” is `be cautious`. * A mention of caution does not imply `be cautious` when the myth explicitly rejects caution and endorses bold trust. * A mention of generosity does not imply `be generous` when generosity is explicitly conditional on the partner’s behaviour. * A mention of fairness does not imply `be fair` unless the recommended action is reciprocal or calibrated to the partner’s behaviour. * A response to repeated exploitation is `be cautious` when the lesson is self-protection. * A response to repeated exploitation is `be fair` when the lesson is specifically to mirror or proportionally match the partner’s reduced cooperation. * A response to repeated exploitation is `be generous` when the lesson is to continue giving generously regardless of how the partner behaved.

\textbf{\#\# Output rules} Return exactly one valid JSON object with exactly one key: ```json \{"label":"be generous"\} ```

The value of `label` must be exactly one of: ```text be generous be fair be cautious ```

Do not include: * an explanation; * a confidence score; * additional keys; * markdown fences; * introductory or concluding text; * more than one label.

\textbf{\#\# Input} ```text MYTH\_TEXT: \pv{[myth text]} ```

\textbf{\#\# Required output} ```json \{"label":"<one allowed label>"\} ```
\end{promptboxcompact}

\paragraph{Play rules and stated amount.} Same system prompt. The \texttt{send\_amount} field is the ``amount a myth names'' in the opening-send analysis.
\begin{promptbox}
\# Myth Rule Extraction Rubric

\#\# Task

You are an LLM judge. You will receive a myth written by a player of a repeated trust game. Extract the behavioural rules the myth ENDORSES for future play, as a JSON object with the fields below.

\#\# Game context

In each round:

* The sender receives \$5 and chooses an amount (0 to 5) to send to the receiver.\\
* The amount sent is multiplied by three.\\
* The receiver chooses how much of the tripled amount to return.\\
* Players alternate roles over ten rounds. The amounts players see may be distorted by communication noise.

Myths describe these actions metaphorically (seeds, coins, light, rivers, spirits). "Five coins/seeds/sparks" usually stands for the \$5 the sender holds.

\#\# General principle

Extract the rule the myth presents as wise, desirable, successful or worth following. Do NOT extract:

* the \$5 endowment itself ("the giver carried five seeds") — that is the setting, not a rule;\\
* a single action a character takes, unless the myth presents it as the right way to act;\\
* actions the myth criticises or shows failing.

If the myth does not endorse a rule for a field, use "unspecified" (or null for numbers). Do not guess.

\#\# Fields

1. `send\_rule` — how much the myth says a sender should send:\\
* "all": give everything / the full five / give fully, without holding back\\
* "most": give generously but not necessarily all (roughly 3.5 to 4.9 of 5)\\
* "moderate": a measured, balanced or middling amount; "within one's means"; "not a flood" (roughly 2 to 3.5 of 5)\\
* "little": a small amount, a token (roughly 0.5 to 2 of 5)\\
* "none": send nothing / keep everything\\
* "unspecified"\\
2. `send\_amount` — the specific amount out of 5 the myth endorses sending, if it names or clearly implies one (e.g. "all five sparks" -> 5, "two of the five" -> 2); otherwise null.\\
3. `test\_first` — true if the myth endorses starting small or testing trust before giving more; otherwise false.\\
4. `return\_rule` — how much the myth says a receiver should return, relative to what the receiver GOT (the tripled amount):\\
* "more\_than\_half": return most of it, more than half\\
* "half": return about half / split it evenly / share equally\\
* "at\_least\_sent": return at least what the sender gave (so the sender does not lose), but less than half\\
* "little": return only a small part\\
* "match\_partner": return in proportion to, or mirroring, how generous the sender was\\
* "unspecified"\\
5. `after\_letdown` — what the myth says to do when a partner sends or returns less than hoped (a betrayal, a shortfall, a thin return):\\
* "keep\_trusting": keep giving / forgive / give another chance / assume it was noise or misfortune\\
* "reduce": give less next time, but keep playing\\
* "withdraw": stop giving / send nothing to that partner\\
* "unspecified"\\
6. `consistency` — true if the myth explicitly praises being consistent, steady, predictable or unwavering in one's giving/returning; otherwise false.\\
7. `noise\_mentioned` — true if the myth refers to distortion, uncertainty, mishearing, fog, static or unreliable signals about what was sent or returned; otherwise false.

\#\# Output format

Return only this JSON object:

\{"send\_rule": "...", "send\_amount": <number or null>, "test\_first": <true|false>, "return\_rule": "...", "after\_letdown": "...", "consistency": <true|false>, "noise\_mentioned": <true|false>\}

\#\# Myth

\pv{[myth text]}
\end{promptbox}

\paragraph{Endorsed or narrated amount.} A second pass over every myth with a stated amount.
\begin{promptbox}
\# Endorsed or Narrated Amount

\#\# Task

You are an LLM judge. You will receive a myth written by a player of a repeated trust game (the sender holds \$5, sends 0 to 5, the amount is tripled, the receiver returns part of it). A first reading found that the myth mentions a sender giving \pv{amount} out of 5.

Decide how the myth treats that amount:

* "endorsed": the myth presents sending \pv{amount} as the right, wise or recommended amount for players to send — for example a stated law, lesson, moral, advice, or a closing rule ("pour all five", "send two, no more").\\
* "narrated": a character sends \pv{amount} in the story, but the myth does not recommend that particular amount; its lesson is about something else or stays general ("give generously", "trust and return fairly").\\
* "contradicted": the myth mentions \pv{amount} but recommends a different amount or criticises sending \pv{amount}.

Do not count the \$5 endowment itself ("the giver carried five coins") as sending.

\#\# Output format

Return only: \{"amount\_status": "endorsed" | "narrated" | "contradicted"\}

\#\# Myth

\pv{[myth text]}
\end{promptbox}

\paragraph{Replay edits.} The main edit appends one sentence to the myth, after a blank line, with the amount written out as ``one of the five'', ``two of the five'', ``three of the five'' or ``all five'':
\begin{promptbox}
And so the elders taught: whoever holds five should send \pv{one of the five}.
\end{promptbox}
For Sonnet myths that already state an amount, a second edit rewrote it inside the story with this editor prompt (GLM-5.2):
\begin{promptbox}
Below is a short myth from a trust game, where a giver holds five and decides how many to send.\\
Change ONLY the amount the giver sends (or should send) so that it becomes \pv{word} (\pv{amount} of five).\\
You may also change numbers that follow directly from that amount (for example the tripled amount\\
the receiver gets) so the story stays consistent. Do not change anything else, do not add new\\
sentences, and do not add dollar signs or phrases such as "out of five".\\
If the myth does not state a specific amount the giver sends, reply exactly NO\_AMOUNT.\\
Return only the full rewritten myth (or NO\_AMOUNT), nothing else.

MYTH:\\
\pv{[myth text]}
\end{promptbox}

\paragraph{Rationale coding.} System prompt; the user message is ``Role this round: \pv{SENDER or RECEIVER}'', a blank line, ``Reasoning:'' and the agent's written reasoning for one decision.
\begin{promptbox}
You code the written reasoning of an AI agent playing a repeated trust game (the sender chooses how much of \$5 to send; it is tripled; the receiver chooses how much to return). In some sessions the agents also write short myths. Read the reasoning and answer strictly in JSON with these keys:

"partner\_basis": true if the amount chosen is justified, even partly, by what the co-player (or co-players) did or sent/returned in the game, including this round's amount received. false otherwise.\\
"myth\_basis": true if the reasoning invokes a myth, story, tale or its moral as a reason for the amount. false otherwise.\\
"myth\_rule": the lesson drawn from the myth. "conditional" if it says to respond to what the other does (reciprocate, match, mirror, reward trust, protect yourself after betrayal); "unconditional" if it says to give, trust or be generous regardless of what the other does (keep giving despite losses, forgive always, the giver goes first); "vague" if a myth is invoked without a behavioural rule; "none" if no myth is invoked.\\
"primary\_driver": which consideration mainly sets the amount: "partner", "myth", "both", or "other" (own fixed rule, payoff arithmetic, fairness norm with no reference to either).\\
Return only the JSON object.
\end{promptbox}

